\section{Players}\label{sec:players}

A seat can be occupied by four kinds of players: a large language model (LLM), a bot, a hybrid, and a decision model. 

\subsection{Language models}

Seats occupied by LLMs are chat models called through OpenRouter. They are given system prompts that only tells them the name of the game, the rules, and the goal (get the highest sum of money), without any strategy advice. In each turn, they are given the full (public) state of the game, consisting of the board, category, used letters, round banks, cards, the wheel, and previous moves. Models play by calling the relevant tool given to them on each turn representing a move. Incorrect tool calls are tried again up to 2 more times before automatically passing.

Each LLM plays with its provider's default temperature and reasoning setting. LLMs are also given their previous reasoning for each move within a round, with prompt caching turned on for all models that support it (Section~\ref{sec:cost}). Appendix~B lists the system prompt and the tool set in its entirely.

We tested seven models: GPT-6.1 Sol and GPT-6.1 Terra (OpenAI), GPT-6 Luna (OpenAI), Claude Haiku 5.5 (Anthropic), DeepSeek V4.1 Flash 1017 (DeepSeek), Qwen 3.8 Flash (Alibaba) and Muse Ember (Meta). Sol and Terra run on OpenAI's half-price flex tier, and we turn on ``adaptive thinking'' for Haiku.

\subsection{The bot}

The bot is a player we built with no LLM component. In this project, we use it as a sort of ``metric'' for the LLMs to play against. 

Bots are equipped with word and phrase lists and frequency derived from wordfreq, Wikipedia, Wiktionary, Wikidata, and Tatoeba's English sentences, with rhyming data from the CMU Pronouncing Dictionary (CMUdict) and category fit from fastText Common Crawl vectors. Given a board, the bot samples 400 answers that fit its requirements and estimates:

\begin{itemize}
  \item each letter's chance of being in the puzzle and expected count of it,
  \item the likeliest answer(s), and
  \item a confidence score tied to that answer based off of how often such answers are correct (e.g. a 70\% confidence in an answer denotes that such answers were right about 70\% of the time, give or take).
\end{itemize}

Bots make their choices with a single value function. In rounds 1 and 2, it scores each available move in expected dollars, aiming to maximize its bank, while in rounds 3 and 4, it aims to finish first. For each move, the value function weighs:

\begin{itemize}
  \item what the move earns immediately, such as the money from a consonant,
  \item what it reveals (e.g. a narrowing down of candidate answers), and
  \item the risk of losing the wheel priced by the chance that an opponent solves before it gets its turn back.
\end{itemize}

The bot estimates that risk for each opponent from how much of the board's letters are visible. During a game, it adjusts the estimate by how often that opponent has solved and acts accordingly (e.g. an opponent that solves early is considered more dangerous and hence solves earlier first to deny them that chance). As vowels cost \$250 and do not yield monetary value, the bot only buys one if the revealing value is worth it.

Once the bot is 95\% or more sure of the answer, a planner takes over. The planner seeks to ``milk'' as much money as possible from the puzzle's leftover consonants, weighed against the chance of Bankrupt/Lose a Turn and the chance of an opponent solving the puzzle before it gets its turn back.

The bot has had six different versions throughout its lifecycle and all comparisons with bots against the other types of players are measured against this version.

\subsection{Hybrids}

Hybrid players take the bot described above and augments it with an LLM assisting its board reading. On each turn, the bot asks the LLM for whether the board is solvable and the answer if so. If the answer fits the board, the bot takes it into account with its own reasoning, with a ``trust curve'' deciding how much the bot factors in the LLM's answer. Despite this, however, the bot is still in full control of its moves. We used GPT-6 Luna and Claude Haiku 5.5 as the LLMs for these hybrid players. 


\subsection{Decision models}

Decision models are also challenged to play this game. Decision models are given an input reflecting legal moves up to this point and return a set of probabilities and a pick, obtained through calling OpenRouter's Decisions API. We always pick what the decision model picks for its gameplay. When a decision model decides to solve, GPT-6 Luna solves the puzzle but otherwise takes no other part in the game in this seat.

We tested d1-ultra (Liquid AI), which ranked first on version 1.0 of Artificial Analysis's Decision Index~\citep{aa-decision-index,liquid2026d1ultra}, and GPT-6 Luna Decisions (OpenAI).
